% generator_I42.tex -- Prop I.42 (Prob.): construct parallelogram equal to given triangle with given angle.
\documentclass[12pt]{article}\usepackage{geometry}\geometry{a5paper,margin=1.5cm}
\usepackage[lining]{ebgaramond}\usepackage{amsmath}\usepackage{amssymb}\usepackage{byrne}
\def\mpPre{u := 1cm; textLabels := false;}
\begin{document}
\defineNewPicture[1/5]{%
  pair A,B,C,D,E,F,G,H,I,J,d[];%
  A:=(0,0);d1:=(8/5u,0);B:=A shifted d1;C:=B shifted (3/2u,0);%
  D:=A shifted (u,-13/5u);E:=D shifted d1;%
  F:=(B--E) intersectionpoint (D--C);%
  G:=2[D,E];d2:=(-xpart(E)+xpart(D)-1/2u,0);%
  H:=B shifted d2;I:=E shifted d2;J:=D shifted d2;%
  draw byPolygon(A,B,F,D)(byyellow);%
  draw byPolygon(D,E,F)(byyellow);%
  draw byPolygon(C,F,E)(byblue);%
  draw byPolygon(E,C,G)(byblack);%
  byAngleDefine(B,E,D,byblue,SOLID_SECTOR);%
  byAngleDefine(H,I,J,byyellow,SOLID_SECTOR);%
  setAttribute("angle","Standalone","HIJ",1);%
  draw byNamedAngleResized();%
  draw byLine(B,E,byred,SOLID_LINE,REGULAR_WIDTH);%
  byLineDefine(A,D,byred,DASHED_LINE,REGULAR_WIDTH);%
  byLineDefine(C,E,byyellow,SOLID_LINE,REGULAR_WIDTH);%
  byLineDefine(A,C,byblue,SOLID_LINE,REGULAR_WIDTH);%
  byLineDefine(D,E,byblack,SOLID_LINE,REGULAR_WIDTH);%
  byLineDefine(E,G,byblack,DASHED_LINE,REGULAR_WIDTH);%
  draw byNamedLineSeq(0)(CE,AC,AD,DE,EG,noLine);%
  draw byLabelsOnPolygon(G,E,D,A,B,C)(ALL_LABELS,0);%
  startAutoLabeling;draw byNamedAngleSidesFull(HIJ)();stopAutoLabeling;%
}
\begin{center}\textbf{Book~I.\enspace Prop.\ XLII. Prob.}\end{center}
\vspace{0.3cm}\noindent
\begin{minipage}[t]{0.50\linewidth}\itshape
To construct a parallelogram equal to a given triangle
\drawPolygon[bottom]{DEF,CFE,ECG} and having an angle equal to a
given rectilinear angle (\drawAngle{I}).
\end{minipage}\hfill
\begin{minipage}[t]{0.46\linewidth}\centering\drawCurrentPicture\end{minipage}
\vspace{0.5cm}
\begin{center}
Make $\drawUnitLine{DE}=\drawUnitLine{EG}$~(pr.~I.10)\\
Draw \drawUnitLine{CE}.\\[4pt]
Make $\drawAngle{E}=\drawAngle{I}$~(pr.~I.23)\\[4pt]
$\left.\begin{aligned}
\mbox{Draw }\drawUnitLine{AD}&\parallel\drawUnitLine{BE}\\
\drawUnitLine{AC}&\parallel\drawUnitLine{DE}
\end{aligned}\right\}$~(pr.~I.31)\\[6pt]
$\drawPolygon{ABFD,DEF}=\mbox{ twice }\drawPolygon{DEF,CFE}$~(pr.~I.41)\\
but $\drawPolygon{DEF,CFE}=\drawPolygon{ECG}$~(pr.~I.38)\\[4pt]
$\therefore \drawPolygon{ABFD,DEF}=\drawPolygon{DEF,CFE,ECG}$.
\end{center}
\begin{flushright}Q.\ E.\ D.\end{flushright}
\end{document}
